In this paper, we introduced relation-centric KGQA, where the goal is to return a subgraph that captures semantic associations among entities, a complementary setting to the standard entity-centric KGQA. 
We presented \textit{UniRel}, a unified modular framework that integrates subgraph retriever with RL-tuned LLMs.
By design, UniRel favors compact and specific subgraphs through a reward function that emphasizes informative relations and low-degree intermediate entities, steering the model toward concise yet distinctive relational answers.
Extensive experiments demonstrate that UniRel consistently improves connectivity and reward over Prompting baselines and generalizes effectively to unseen entities and relations.
Our analysis further revealed a model-level distinction: Qwen models exhibit a stronger reliance on semantic cues, whereas Llama models maintain more stable performance by leveraging structural connectivity.
Also, our results on multi-entity queries confirm the scalability of UniRel, showing that the framework extends naturally to complex reasoning tasks.
Beyond the relation-centric setting, we further showed that the UniRel pipeline can be applied to conventional entity-centric KGQA by replacing the subgraph retriever module, achieving competitive or improved Hit@1 and F1 scores across multiple benchmarks.




% In this paper, we introduced relation-centric KGQA, where the goal is to return a subgraph that captures semantic associations among entities, a complementary setting to the standard entity-centric KGQA. 
% We presented \textit{UniRel-R1}, a unified framework that integrates subgraph selection, multi-stage pruning, and RL-tuned LLMs, where the reward design explicitly favors compact subgraphs with rarer relations and lower-degree intermediates, thereby steering the model toward concise and distinctive relational answers.
% Through extensive evaluations on seven benchmark knowledge graphs, we demonstrated that UniRel-R1 delivers substantial improvements over Prompting baselines, achieving large gains in both connectivity and reward while maintaining strong generalization to unseen entities and relations.
% Our analysis further revealed a model-level distinction: Qwen models exhibit a stronger reliance on semantic cues, whereas Llama models maintain more stable performance by leveraging structural connectivity.
% Finally, our results on multi-entity queries confirm the scalability of UniRel-R1, showing that the framework extends naturally to complex reasoning tasks.